\section{随机过程}
\subsection{随机过程的基本概念}
\subsubsection{定义}
\textbf{随机函数}
\begin{itemize}
    \item 随机变量的集合;
    \item 所有样本函数的集合.
\end{itemize}

\subsubsection{分布函数}
\textbf{一维分布函数}
\begin{equation}
    F_X(x;t)=P[X(t)\leq x].
\end{equation}

\textbf{一维概率密度函数}
\begin{equation}
    f_X(x;t)=\frac{\partial F_X(x;t)}{\partial x}.
\end{equation}

\subsubsection{数字特征}
\textbf{数学期望}
\begin{equation}
    a(t)=E[X(t)]=\int_{-\infty}^{\infty}xf_X(x;t)\mathrm{d}x.
\end{equation}

\textbf{方差}
\begin{equation}
    \sigma^2(t)=D[X(t)]=E[X^2(t)]-a^2(t).
\end{equation}

\textbf{自相关函数}
\begin{equation}
    \begin{aligned}
        R(t,t+\tau) & =E[X(t)X(t+\tau)]                                                                                     \\
                    & =\int_{-\infty}^{\infty}\int_{-\infty}^{\infty}x_1x_2f_X(x_1,x_2;t,t+\tau)\mathrm{d}x_1\mathrm{d}x_2.
    \end{aligned}
\end{equation}

\textbf{互相关函数}
\begin{equation}
    R_{XY}(t,t+\tau)=E[X(t)Y(t+\tau)].
\end{equation}
